% !PRINTED 98
\weberpage{98}
%<<CONT
neation $\mathfrak{M}$ die Rolle der Einheit. Ist $S = \begin{pmatrix} a, & b \\ c, & \partial \end{pmatrix}$ eine beliebige Transformation, so ist
\[
\begin{pmatrix} a, & b \\ c, & \partial \end{pmatrix}\begin{pmatrix} \partial, & -b \\ -c, & a \end{pmatrix} = \begin{pmatrix} n, & 0 \\ 0, & n \end{pmatrix}
\]
eine Multiplikation. Die beiden Kollineationen
\[
C = \mathfrak{M}\begin{pmatrix} a, & b \\ c, & \partial \end{pmatrix}, \quad C^{-1} = \mathfrak{M}\begin{pmatrix} \partial, & -b \\ -c, & a \end{pmatrix}
\]
geben also bei der Komposition $CC^{-1} = C^{-1}C = \mathfrak{M}$ und sind also zueinander reziprok.

\begin{center}
\gesperrt{Demnach bildet die Gesamtheit der Kollineationen eine Gruppe.}
\end{center}

Die Transformationen vom Grade 1 heißen \gesperrt{lineare Transformationen}. Wir bezeichnen bei diesen die Transformationszahlen zum Unterschiede mit den griechischen Buchstaben $\alpha, \beta, \gamma, \delta$, so daß
\[
L = \begin{pmatrix} \alpha, & \beta \\ \gamma, & \delta \end{pmatrix}, \quad \alpha\delta - \beta\gamma = 1
\]
eine lineare Transformation bedeutet. Das System $\mathfrak{L}$ der linearen Transformationen ist eine in $\mathfrak{S}$ enthaltene Gruppe, denn sind
\[
L = \begin{pmatrix} \alpha, & \beta \\ \gamma, & \delta \end{pmatrix}, \quad L' = \begin{pmatrix} \alpha', & \beta' \\ \gamma', & \delta' \end{pmatrix},
\]
so ist
\[
L'' = L'L = \begin{pmatrix} \alpha'\alpha+\beta'\gamma, & \alpha'\beta+\beta'\delta \\ \gamma'\alpha+\delta'\gamma, & \gamma'\beta+\delta'\delta \end{pmatrix} = \begin{pmatrix} \alpha'', & \beta'' \\ \gamma'', & \delta'' \end{pmatrix}
\]
gleichfalls linear, und man kann $L$ bei gegebenem $L'$, $L''$ aus den Gleichungen
\[
\begin{aligned}
\alpha'\alpha+\beta'\gamma &= \alpha'', & \alpha'\beta+\beta'\delta &= \beta'', \\
\gamma'\alpha+\delta'\gamma &= \gamma'', & \gamma'\beta+\delta'\delta &= \delta''
\end{aligned}
\]
eindeutig bestimmen. Die Einheit der Gruppe $\mathfrak{L}$ ist die identische Substitution $\begin{pmatrix} 1, & 0 \\ 0, & 1 \end{pmatrix}$ und jede Substitution $L$ hat ihre Reziproke $L^{-1}$, wie aus der Zusammensetzung
\[
LL^{-1} = \begin{pmatrix} \alpha, & \beta \\ \gamma, & \delta \end{pmatrix}\begin{pmatrix} \delta, & -\beta \\ -\gamma, & \alpha \end{pmatrix} = \begin{pmatrix} 1, & 0 \\ 0, & 1 \end{pmatrix}
\]
hervorgeht. Die Gruppe $\mathfrak{L}$ ist unendlich und ist \gesperrt{nicht kommutativ}.
